% ══════════════════════════════════════════════════════
% Giao diện theo mẫu slide môn học (CS313 — Course Overview):
% Beamer Madrid + colortheme beaver (đỏ đậm trên nền xám), font Latin Modern Sans
% (bản OpenType của Computer Modern Sans, có đủ dấu tiếng Việt). Biên dịch bằng XeLaTeX.
% ══════════════════════════════════════════════════════

\usetheme{Madrid}
\usecolortheme{beaver}
\usepackage{tikz}   % theme cũ (Durham) tự nạp TikZ; phần 1 có dùng

% ── Font: Latin Modern Sans (nạp theo tên file để chạy được cả trên Overleaf lẫn máy cá nhân) ──
\usepackage{fontspec}
\setsansfont{lmsans10-regular.otf}[
  BoldFont       = lmsans10-bold.otf,
  ItalicFont     = lmsans10-oblique.otf,
  BoldItalicFont = lmsans10-boldoblique.otf]
\setmonofont{lmmono10-regular.otf}

% ── Tên màu cũ vẫn được các phần slide dùng → ánh xạ sang bảng màu mới ──
\colorlet{Theme_Color}{darkred}             % màu nhấn chính (đỏ của beaver)
\colorlet{RiverDark}{darkred!80!black}      % đỏ đậm
\colorlet{RiverLight}{gray!10}              % nền xám nhạt (giống nền tiêu đề slide)
\colorlet{RiverAccent}{darkred!50}

% ── Block: Madrid mặc định dùng tiêu đề xanh navy → đổi sang tông đỏ/xám cho khớp beaver ──
\setbeamercolor{block title}{fg=white, bg=darkred!80!black}
\setbeamercolor{block title alerted}{fg=white, bg=darkred}
\setbeamercolor{block title example}{fg=white, bg=gray!60!black}
\setbeamercolor{block body}{bg=gray!8}
\setbeamercolor{block body alerted}{bg=gray!8}
\setbeamercolor{block body example}{bg=gray!8}

% ── Trang phân mục & các slide dùng chung ──
\renewcommand{\sectionname}{Phần}
\AtBeginSection[]{\begin{frame}\sectionpage\end{frame}}
\newcommand{\makeoutline}{\begin{frame}{Nội dung}\tableofcontents\end{frame}}
\newcommand{\makethankyou}{%
  \begin{frame}
    \centering\vfill
    {\usebeamercolor[fg]{title}\Large Cảm ơn thầy cô và các bạn!}\\[1em]
    Câu hỏi · Góp ý · Thảo luận
    \vfill
  \end{frame}}
